
\documentclass[12pt,a4paper]{article}

\usepackage[utf8]{inputenc}
\usepackage[T1]{fontenc}
\usepackage[brazil]{babel}
\usepackage{amsmath}
\usepackage{siunitx}
\usepackage{booktabs}
\usepackage{graphicx}
\usepackage[margin=2.5cm]{geometry}

\graphicspath{{figs/}}

\begin{document}


\subsection*{Caso 2J}
\label{sec:res-caso2.2}
No \texttt{Caso 2J} Substitui-se a hipótese de um tubo perfeitamente retilíneo pela tortuosidade anatômica real do NO, imposta diretamente à linha de centro do complexo dural-neural na forma de uma curvatura em ``J''. Preservadas as propriedades de malha, materiais e condições de contorno, o \texttt{Caso 2J} difere do \texttt{Caso 2G} em única variável da geometria curva, isolando de modo controlado o efeito mecânico da tortuosidade.

O primeiro efeito mecânico mensurável da tortuosidade é o amolecimento axial inicial do conjunto. A rigidez axial secante pré-flambagem (recuo $\approx \SI{0.10}{\milli\meter}$, ancorada na origem), medida pela componente de reação ao longo do eixo de compressão do nervo, cai de $k_{\mathrm{axial}} \approx \SI{306}{\newton\per\meter}$ (tubo reto) para $k_{\mathrm{axial}} \approx \SI{178}{\newton\per\meter}$ (curvatura em J), uma redução de aproximadamente \SI{42}{\%}; ao final da rampa clínica ambos enrijecem geometricamente. O deslocamento lateral máximo do nervo atinge $\approx \SI{0.5}{\milli\meter}$ para um recuo axial de \SI{1.0}{\milli\meter}.

No entanto, a razão entre os dobramentos da pia e da dura permanece $\approx 1{,}0$, indicando que pia e dura fletem solidariamente, exatamente o regime que a teoria de Hetényi prevê para a fundação elástica de uma órbita saudável ($k_w = \SI{200}{\kilo\pascal\per\meter}$). Isso revela que o \texttt{Caso 2J} produz uma flexão global de todo o feixe em um único arco, e não a ondulação focal de duas meias-ondas antissimétricas (modo $n = 2$), como padrão antissimétrico em ``S'' clinicamente observado na SANS. 

Além disso, para isolar o papel da gordura orbital, variou-se a rigidez da fundação de Winkler por um fator de $250\times$ (cerca de duas ordens e meia de grandeza, de $k_w = \SI{20}{\kilo\pascal\per\meter}$ a $\SI{5}{\mega\pascal\per\meter}$), mantendo todo o restante fixo (Figura~\ref{fig:res-caso2j-winkler}). Avaliando no limite clínico de recuo axial de \SI{1.0}{\milli\meter}, enrijecer a gordura apenas \emph{confina} a tortuosidade, reduzindo a excursão lateral da pia de \SI{0.65}{\milli\meter} (gordura jovem macia) para \SI{0.065}{\milli\meter} (gordura SANS-like, $25\times$ mais rígida) e a reação no engaste de \SI{485}{} para \SI{342}{\milli\newton}, sem jamais converter o arco em uma ondulação antissimétrica em ``S''.

\begin{figure}[ht]
\centering
\includegraphics[width=0.8\textwidth]{Figuras/Mestrado/on-caso-2j-winkler.png}
\caption{Parametrização da gordura orbital (fundação de Winkler) no \texttt{Caso 2J}, avaliada no limite clínico de recuo axial de \SI{1.0}{\milli\meter}. Ao enrijecer a gordura de \SI{20}{\kilo\pascal\per\meter} a \SI{5}{\mega\pascal\per\meter}, tanto a excursão lateral máxima da pia (eixo esquerdo) quanto a reação no engaste posterior (eixo direito) caem monotonicamente; em todo o intervalo a coluna atinge $\lambda = 1{,}0$ sem instabilidade, mantendo o arco difuso de modo $n = 1$.}
\label{fig:res-caso2j-winkler}
\end{figure}

\end{document}
